\chapter{How a program becomes a running thing}
\label{ch:01}

You wrote a CPU. In Book~5 you designed an 8-bit ALU, wired it into a
datapath, and drove it with a two-state FSM until Fibonacci numbers came
out of a serial pin. That machine is small enough to hold in your head ---
302 synthesis cells, ten opcodes, two states --- but it is not a toy. The
CPU you are using right now runs \emph{more} code than yours, not
\emph{different} code. Somewhere between \texttt{int main(void)} and the
transistors of your cpu8 there is a journey every program takes, and this
book is about travelling it on foot. This chapter walks the first mile: how
source text becomes a file of machine code, how that file gets into memory,
and where your cpu8 meets the real silicon. Along the way you will
compile, disassemble and link \emph{real} binaries and read their real
output --- nothing in this chapter is simulated.

\index{compilation pipeline}\index{ELF}\index{linker}\index{loader}%
\index{relocation}\index{symbol table}\index{object file}%
\index{execve@\texttt{execve}}

\section{The four-stage journey}

\index{preprocessor}\index{compiler}\index{assembler}
Type \texttt{cc sum.c -o sum} and four programs run in sequence, each
handing a new file to the next. The names differ between toolchains ---
\texttt{gcc}, \texttt{clang}, \texttt{cc} --- the shape does not:

\begin{enumerate}
  \item \textbf{Preprocess} --- text in, text out. \texttt{\#include}
        splices header files in place; \texttt{\#define} substitutes
        macros; conditionals (\texttt{\#ifdef}) delete code before the
        compiler ever sees it. Output: a ``translation unit'', still plain
        C.
  \item \textbf{Compile} --- the translation unit becomes assembly:
        mnemonics, register names, labels --- the same language your cpu8
        assembler spoke, with a bigger vocabulary. Output: \texttt{.s}.
  \item \textbf{Assemble} --- mnemonics become machine-code bytes plus
        bookkeeping (which bytes are code, which are data, which names are
        exported). Output: an \emph{object file}, \texttt{.o} --- already
        ELF, just not yet executable.
  \item \textbf{Link} --- object files are stitched together: symbols
        resolved, sections merged, addresses assigned, one entry point
        recorded. Output: an executable ELF.
\end{enumerate}

\begin{example}[The pipeline, hands on]
\label{ex:pipeline}
Here is the whole journey on \texttt{code/ch01\_sum.c}, a six-line
freestanding sum. Each stage was run on the dev box
(Apple clang 21, targeting x86-64 Linux so the output is real ELF):

\begin{lstlisting}[basicstyle=\ttfamily\scriptsize]
clang -E     ch01_sum.c > ch01_sum.i     # 1: preprocess
clang -S     ch01_sum.i                  # 2: compile -> ch01_sum.s
clang -c     ch01_sum.s                  # 3: assemble -> ch01_sum.o
ld.lld -o ch01_sum.elf ch01_start.o ch01_sum.o   # 4: link
\end{lstlisting}

Stage 2 emits assembly you can already read. Real output from
\texttt{clang -S} (first lines of \texttt{sum}):

\begin{lstlisting}[basicstyle=\ttfamily\scriptsize]
sum:                                    # @sum
        pushq   %rbp
        movq    %rsp, %rbp
        movq    %rdi, -8(%rbp)
        movl    %esi, -12(%rbp)
        movl    $0, -16(%rbp)           # s = 0
        movl    $0, -20(%rbp)           # i = 0
.LBB0_1:                                # loop:
        movl    -20(%rbp), %eax         #   eax = i
        cmpl    -12(%rbp), %eax         #   compare i, n
        jge     .LBB0_4                 #   done when i >= n
        movq    -8(%rbp), %rax
        movslq  -20(%rbp), %rcx
        movl    (%rax,%rcx,4), %eax     #   eax = a[i]
        addl    %eax, -16(%rbp)         #   s += a[i]
\end{lstlisting}

Every line is one instruction the CPU will execute; \texttt{-8(\%rbp)}
is the array pointer spilled to the stack exactly the way your cpu8's
indirect addressing would have handled it.
\end{example}

\begin{figure}[htbp]
\centering
\begin{tikzpicture}[font=\scriptsize, >=latex]
% the four-stage journey
\foreach \i/\t/\d/\o in {
  0/{preprocess}/{\#include, \#define}/{.c},
  1/{compile}/{parse, optimise, emit asm}/{.s},
  2/{assemble}/{mnemonics $\to$ bytes}/{.o},
  3/{link}/{resolve symbols, lay out}/{ELF exec}}{
  \pgfmathsetmacro{\x}{0.2+\i*2.5}
  \draw[thick, fill=accentsoft] (\x,3.0) rectangle (\x+2.2,4.3);
  \node[align=center] at (\x+1.1,3.9) {\t};
  \node[align=center, font=\tiny, text=slate] at (\x+1.1,3.35) {\d\\\ttfamily \o};
}
\foreach \i in {0,...,2}{
  \pgfmathsetmacro{\x}{2.4+\i*2.5}
  \draw[->, thick] (\x,3.65) -- (\x+0.3,3.65);}
% load stage
\draw[thick, fill=goodbg] (8.2,1.0) rectangle (10.4,2.3);
\node[align=center] at (9.3,1.9) {load (\texttt{execve})};
\node[align=center, font=\tiny, text=slate] at (9.3,1.4) {map text/data,\\jump to entry};
\draw[->, thick] (9.3,3.0) -- (9.3,2.3);
\node[align=left, text=slate, text width=10.3cm] at (5.2,0.2)
{Four tool invocations, one machine read of the result. The .o is
already ELF --- linking only decides where each byte goes; loading
only decides which page it lands on.};
\end{tikzpicture}

\caption{The four-stage journey from source to running program. Each
stage's product is a file you can inspect; none is magic.}
\label{fig:pipeline01}
\end{figure}

\begin{checkpoint}
In one sentence each: what does the preprocessor hand to the compiler,
and what does the assembler hand to the linker?
\end{checkpoint}
\begin{solution}
The preprocessor hands \emph{expanded C source} (one translation unit,
no \texttt{\#}-lines) to the compiler. The assembler hands the linker an
\emph{object file}: machine-code bytes plus a symbol table and
relocation records --- not yet a program, but already ELF.
\end{solution}

\begin{tripwire}[The pipeline has four stages, not one]
When a student writes ``gcc compiles my program'', they hide three of the
four jobs under one word. You will meet students who cannot tell a
\emph{compile} error (bad syntax) from a \emph{link} error (missing
symbol) --- the fix is different every time, and the fix-time is the
difference between seconds and an afternoon.
\end{tripwire}

\section{Object files and ELF sections}
\label{sec:elf01}

\index{section (ELF)}\index{.text@\texttt{.text}}\index{.data@\texttt{.data}}
An object file is not a blob: it is a filing cabinet. The bytes from
stage~3 arrive sorted into \emph{sections}, each with a name, a size, a
job. Compile the sum program and look inside --- real output, dev box,
clang 21 targeting x86-64 ELF:

\begin{lstlisting}[basicstyle=\ttfamily\scriptsize]
$ objdump -h ch01_sum.o
Sections:
Idx Name            Size     VMA              Type
  1 .strtab         00000073 0000000000000000
  2 .text           00000076 0000000000000000 TEXT
  3 .rela.text      00000030 0000000000000000
  4 .data           00000014 0000000000000000 DATA
  7 .eh_frame       00000058 0000000000000000 DATA
  8 .rela.eh_frame  00000030 0000000000000000
\end{lstlisting}

Read the four that matter:

\begin{itemize}
  \item \texttt{.text} --- 0x76 = 118 bytes of machine code. This is what
        a CPU's instruction fetch unit reads. It is called ``text'' for
        the same reason a teletype's paper tape was: the name is older
        than most programmers, and it means \emph{code}, not English.
  \item \texttt{.data} --- 0x14 = 20 bytes of initialised data: the
        five-element array \texttt{v} inside \texttt{main}, stored as
        real bytes.
  \item \texttt{.rela.text} --- 48 bytes of relocation records: places in
        \texttt{.text} the linker must patch once it knows the final
        addresses.
  \item \texttt{.symtab} (hidden in this excerpt, present in the file)
        --- the symbol table: the names \texttt{sum}, \texttt{main},
        whether each is defined or referenced, and where.
\end{itemize}

\begin{example}[Find \texttt{main}'s machine code]
\label{ex:disasm}
\texttt{objdump -d} walks the \texttt{.text} section and prints every
instruction with its offset and its encoding bytes. Real output for
\texttt{main} in \texttt{ch01\_sum.o}:

\begin{lstlisting}[basicstyle=\ttfamily\scriptsize]
0000000000000050 <main>:
  50: 55                    pushq   %rbp
  51: 48 89 e5              movq    %rsp, %rbp
  54: 48 83 ec 10           subq    $0x10, %rsp
  58: c7 45 fc 00 00 00 00  movl    $0x0, -0x4(%rbp)
  5f: 48 8d 3d 00 00 00 00  leaq    (%rip), %rdi  # 0x66
  66: be 05 00 00 00        movl    $0x5, %esi
  6b: e8 00 00 00 00        callq   0x70 <main+0x20>
  70: 48 83 c4 10           addq    $0x10, %rsp
  74: 5d                    popq    %rbp
  75: c3                    retq
\end{lstlisting}

Three things to notice. First, the middle column is the machine code ---
\texttt{55}, \texttt{48 89 e5}: the exact bytes Book~5 ch16's ROM would
have held. Second, \texttt{callq} at \texttt{6b} encodes
\texttt{00 00 00 00} for its target: the assembler left a hole because it
did not yet know where \texttt{sum} would live --- the linker's relocation
step fills it. Third, \texttt{movl \$0x5, \%esi} is your accumulator
pattern: load a constant into a register, then hand it to the call.
\end{example}

\begin{figure}[htbp]
\centering
\begin{tikzpicture}[font=\scriptsize, >=latex]
% ELF anatomy of the real ch01_sum.elf
\draw[thick, fill=graybg] (0.6,0.6) rectangle (4.4,4.6);
\node[font=\footnotesize\bfseries] at (2.5,4.3) {ELF header};
\foreach \y/\t/\s in {
  3.8/{.text (0x96 B)}/{the code, vma 0x2011c0},
  3.2/{.data (0x34 B)}/{initialised data},
  2.6/{.eh\_frame}/{unwind info},
  2.0/{.symtab + .strtab}/{names for linkers},
  1.4/{.shstrtab + .comment}/{section names, meta}}{
  \draw[thick] (0.8,\y-0.24) rectangle (4.2,\y+0.24);
  \node[align=left, font=\tiny] at (0.95,\y) {\t};
}
\node[font=\tiny, text=slate] at (2.5,0.9) {real \ttfamily ch01\_sum.elf};
% section headers annotation
\draw[->, thick, slate] (4.4,3.8) -- (6.0,3.8)
  node[right, align=left, font=\tiny, text width=4.2cm]
  {code bytes, executable --- the part your cpu8 would fetch};
\draw[->, thick, slate] (4.4,2.0) -- (6.0,2.0)
  node[right, align=left, font=\tiny, text width=4.2cm]
  {symbol table --- the linker's contract, invisible at run time};
\node[align=left, text=slate, text width=10cm] at (5.3,-0.2)
{An ELF is a table of contents plus the bytes each section points
at: same bytes, four different jobs.};
\end{tikzpicture}

\caption{The section map of the real \texttt{ch01\_sum.elf}: code, data,
bookkeeping --- one file, four jobs.}
\label{fig:elf01}
\end{figure}

\begin{tripwire}[\texttt{.text} is not text, \texttt{.data} is not ``all
the data'']
New readers lose hours to naming. \texttt{.text} = instructions.
\texttt{.data} = only \emph{initialised writable} globals.
\texttt{.bss} = initialised-to-zero globals (not stored in the file at
all --- the loader zero-fills it, which is why a 1~MB zero array costs
no disk space).
\end{tripwire}

\section{Symbols and relocation}
\label{sec:sym01}

\index{symbol (linker)}\index{relocation record}\index{undefined symbol}
The symbol table is the linker's contract. Two kinds of entries decide
whether a program links: \emph{defined} symbols (the file provides them)
and \emph{undefined} symbols (the file promises to get them from
somewhere else). \texttt{objdump -t} shows both:

\begin{lstlisting}[basicstyle=\ttfamily\scriptsize]
$ objdump -t ch01_sum.elf          # after linking
00000000002011c0 g       .text   _start
0000000000202260 g       .data   data
00000000002011e0 g     F .text   sum          # defined function
0000000000201230 g     F .text   main         # defined function
\end{lstlisting}

While the object is still a \texttt{.o}, the assembler records a
\emph{relocation} for every address it could not finish. Real
\texttt{objdump -r} on \texttt{ch01\_broken.o}:

\begin{lstlisting}[basicstyle=\ttfamily\scriptsize]
RELOCATION RECORDS FOR [.text]:
OFFSET           TYPE            VALUE
0000000000000012 R_X86_64_PC32   .data-0x4
000000000000001c R_X86_64_PLT32  sum-0x4
0000000000000024 R_X86_64_PLT32  missing-0x4
\end{lstlisting}

Each row says: ``patch \texttt{.text} at this offset, against the
final address of that symbol.'' The linker runs through these records,
looks up every undefined symbol in every other object's symbol table,
and either patches the hole --- or reports the failure you saw live:

\begin{lstlisting}[basicstyle=\ttfamily\scriptsize]
$ ld.lld -o broken ch01_broken.o
ld.lld: error: undefined symbol: sum
>>> referenced by ch01_broken.c
>>>            ch01_broken.o:(main)
ld.lld: error: undefined symbol: missing
>>> referenced by ch01_broken.c
>>>            ch01_broken.o:(main)
\end{lstlisting}

\begin{figure}[htbp]
\centering
\begin{tikzpicture}[font=\scriptsize, >=latex]
% symbol + relocation
\draw[thick, fill=accentsoft] (0.2,2.8) rectangle (4.4,4.6);
\node[align=center, font=\footnotesize] at (2.3,4.1) {sum.o};
\node[align=left, font=\tiny, text width=3.9cm] at (2.3,3.4)
{\texttt{F sum} defined here\\symbol table: \texttt{0000... g F .text sum}};
\draw[thick, fill=graybg] (5.4,2.8) rectangle (9.6,4.6);
\node[align=center, font=\footnotesize] at (7.5,4.1) {broken.o};
\node[align=left, font=\tiny, text width=3.9cm] at (7.5,3.4)
{\texttt{U sum}, \texttt{U missing}\\undefined --- the call sites hold
a promise, not an address};
\draw[->, thick] (4.4,3.7) -- (5.4,3.7)
  node[midway, above, font=\tiny] {link};
\draw[thick, fill=goodbg] (3.6,0.8) rectangle (6.2,2.2);
\node[align=center, font=\tiny] at (4.9,1.7) {resolved: \texttt{sum}};
\node[align=center, font=\tiny, text=amber] at (4.9,1.25) {unresolved:
\texttt{missing} $\to$ error};
\draw[->, thick] (7.5,2.8) -- (6.0,1.7);
\node[align=left, text=slate, text width=10cm] at (5.1,-0.4)
{Relocation = the linker's bookkeeping: every undefined \texttt{U}
must find a defined \texttt{F/g} in another object --- or the real
error you just saw appears.};
\end{tikzpicture}

\caption{Relocation is bookkeeping: the assembler records holes, the
linker resolves them --- or names the missing symbol.}
\label{fig:reloc01}
\end{figure}

\begin{example}[Hand-link, and watch the error message change]
\label{ex:handlink}
Link \texttt{ch01\_sum.o} alone and \texttt{ld.lld} reports
\texttt{undefined symbol: sum} plus \texttt{undefined symbol: missing} ---
two separate records, two separate complaints. Add \texttt{ch01\_sum.o}
to the link line and \texttt{sum} disappears from the error list but
\texttt{missing} remains. That is the whole game: undefined symbols are
holes to fill; a link is complete when every hole has a defined home.
\end{example}

\begin{checkpoint}
Your build reports \texttt{undefined symbol: \_\_stack\_chk\_fail}.
Is this a compile error or a link error, and what does the linker want?
\end{checkpoint}
\begin{solution}
A link error --- the compiler emitted a call to a helper that exists in
the C library but is not on your link line. The linker wants any object
file or library that \emph{defines} that symbol; adding it is the fix,
not editing your source.
\end{solution}

\section{The loader and \texttt{execve}}
\label{sec:exec01}

\index{execve@\texttt{execve}}\index{address space}\index{page table}
Linking produced a file; something still has to make it \emph{run}. That
is the loader's job, and on Linux it is invoked by \texttt{execve}: the
kernel opens the ELF, checks the magic bytes \texttt{\textbackslash
x7f~E~L~F}, builds a fresh address space, maps the loadable segments to
virtual pages, sets the program counter to the entry point, and hands
the CPU to it. After the link above, the real executable's sections
look like this:

\begin{lstlisting}[basicstyle=\ttfamily\scriptsize]
$ objdump -h ch01_sum.elf
Sections:
Idx Name          Size     VMA              Type
  1 .eh_frame     0000005c 0000000000200158 DATA
  2 .text         00000096 00000000002011c0 TEXT
  3 .data         00000034 0000000000202260 DATA
  5 .symtab       000000a8 0000000000000000
\end{lstlisting}

Compare with the \texttt{.o}: the \texttt{VMA} column is no longer zero.
The linker assigned real addresses --- \texttt{.text} at
\texttt{0x2011c0} --- and the loader maps those sections at exactly those
virtual addresses. The program counter lands at \texttt{\_start}
(\texttt{0x2011c0}), and the machine runs.

\begin{figure}[htbp]
\centering
\begin{tikzpicture}[font=\scriptsize, >=latex]
% process address space, low->high
\draw[thick] (0.6,0.4) rectangle (2.4,4.8);
\foreach \y/\t/\c in {
  4.4/{high: stack (grows down)}/{graybg},
  3.5/{mmap region}/{white},
  2.7/{heap (grows up)}/{graybg},
  2.1/{.bss + .data}/{white},
  1.5/{.text}/{accentsoft},
  0.8/{low: unused + kernel-reserved}/{white}}{
  \draw[fill=\c] (0.6,\y-0.35) rectangle (2.4,\y+0.35);
  \node[font=\tiny] at (1.5,\y) {\t};
}
\node[right, align=left, font=\tiny, text width=6.8cm] at (3.0,3.5)
{the ELF's PT\_LOAD segments become text+data here};
\node[right, align=left, font=\tiny, text width=6.8cm] at (3.0,2.7)
{the loader builds this space; the CPU only sees the addresses};
\node[right, align=left, font=\tiny, text width=6.8cm] at (3.0,1.5)
{entry point jumps here --- your cpu8's pc=0x0, grown up};
\node[align=left, text=slate, text width=10cm] at (5.3,-0.4)
{Same layout idea as the machine you built in Book 5, same reason:
the hardware reads addresses, the software draws the map.};
\end{tikzpicture}

\caption{The process address space the loader builds: the ELF's
segments become mapped pages; the CPU reads only addresses.}
\label{fig:address01}
\end{figure}

\begin{example}[Where the map meets the metal]
\label{ex:bridge}
On your cpu8, \texttt{pc=0x0} meant ROM offset zero. On a real machine,
\texttt{pc=0x2011c0} means ``page \texttt{0x20000} plus \texttt{0x1c0}'',
and a page table translates that to DRAM. The idea is identical ---
\emph{the CPU fetches from an address, and something decides what lives
there}. What changed is who decides: on your cpu8, you did; on Linux,
the page table your OS maintains does. That is the one sentence this
book is about.
\end{example}

\section{Where Book 5's CPU meets real silicon}
\label{sec:bridge01}

\index{RV32I}\index{cpu8}\index{Book 5}
Book~5's cpu8 and a RISC-V RV32I are the same machine at different
scales. The mapping table you already know:

\begin{center}
\begin{tabularx}{\linewidth}{lX}
\toprule
\textbf{cpu8 (Book 5 ch14--16)} & \textbf{RV32I (this book)} \\
\midrule
8-bit regs, r0 = accumulator A &
   32 32-bit regs, x0 wired to 0 \\
10 opcodes, 4-bit op field &
   RV32I base: 37 instructions, 7-bit opcode \\
2-state FSM: fetch, execute &
   5-stage pipeline: IF ID EX MEM WB \\
\texttt{JMP}/\texttt{JZ} flags &
   \texttt{beq}/\texttt{bne}/\dots compare-and-branch \\
16-byte ROM, \texttt{pc=0x0} &
   virtual memory, \texttt{pc} = entry VA \\
UART prints a byte &
   \texttt{ecall} asks the OS to print \\
\bottomrule
\end{tabularx}
\end{center}

The right column is not a different kind of computer --- it is the same
datapath+FSM shape with more bits and more stages. Chapters 5--8 build
the pipelined RV32I core for real; this section is the signpost that
says you have already built its skeleton once.

\begin{figure}[htbp]
\centering
\begin{tikzpicture}[font=\scriptsize, >=latex]
% book5 cpu8 -> real machine hand-off
\draw[thick, fill=accentsoft] (0.3,1.4) rectangle (4.1,3.6);
\node[align=center, font=\footnotesize] at (2.2,3.1) {Book 5, ch14--16};
\node[align=left, font=\tiny, text width=3.5cm] at (2.2,2.3)
{cpu8: 302 cells, 10 opcodes,\\
2-state FSM, Fibonacci\\on 12 real OUT writes};
\draw[thick, fill=goodbg] (5.9,1.4) rectangle (9.7,3.6);
\node[align=center, font=\footnotesize] at (7.8,3.1) {Book 7, ch5--8};
\node[align=left, font=\tiny, text width=3.5cm] at (7.8,2.3)
{pipelined RV32I: 5 stages,\\
forwarding, hazards, traps ---\\
the same shape, real speed};
\draw[->, very thick, accent] (4.1,2.5) -- (5.9,2.5)
  node[midway, above, font=\tiny] {this book};
\node[align=left, text=slate, text width=10cm] at (5.0,0.5)
{The hand-off this book exists for: the machine you built proves the
shape; the chapters ahead prove it at scale.};
\end{tikzpicture}

\caption{The hand-off: Book 5 proved the shape; this book proves it at
scale.}
\label{fig:handoff01}
\end{figure}

\section{A map of the book}
\label{sec:map01}

Six parts, one question each. The order matters: each part runs the
tools the previous one installed.

\begin{figure}[htbp]
\centering
\begin{tikzpicture}[font=\scriptsize, >=latex]
% book map
\foreach \x/\t/\d in {
  0.2/{Part I (ch1--4)}/{source $\to$ machine code},
  3.4/{Part II (ch5--8)}/{the RV32I capstone},
  6.6/{Part III (ch9--11)}/{memory hierarchy}}{
  \draw[thick, fill=accentsoft] (\x,3.4) rectangle (\x+2.9,4.8);
  \node[align=center, font=\footnotesize] at (\x+1.45,4.4) {\t};
  \node[align=center, font=\tiny] at (\x+1.45,3.9) {\d};
}
\foreach \x/\t/\d in {
  1.8/{Part IV (ch12--14)}/{the OS: processes, files},
  5.0/{Part V (ch15--17)}/{performance craft},
  8.2/{Part VI (ch18)}/{runtimes}}{
  \draw[thick, fill=graybg] (\x,1.6) rectangle (\x+2.9,3.0);
  \node[align=center, font=\footnotesize] at (\x+1.45,2.6) {\t};
  \node[align=center, font=\tiny] at (\x+1.45,2.1) {\d};
}
\draw[->, thick] (3.1,4.1) -- (3.4,4.1);
\draw[->, thick] (6.3,4.1) -- (6.6,4.1);
\draw[->, thick] (7.9,3.4) -- (6.8,3.0);
\draw[->, thick] (4.7,3.0) -- (4.0,3.4);
\draw[->, thick] (6.5,3.0) -- (6.6,3.4);
\node[align=left, text=slate, text width=10cm] at (5.2,0.6)
{The path this chapter starts: build the real machine around the
toy CPU you already made, then measure what the folklore got
wrong.};
\end{tikzpicture}

\caption{The roadmap: source $\to$ machine code $\to$ a real core
$\to$ memory $\to$ the OS $\to$ craft.}
\label{fig:roadmap01}
\end{figure}

\begin{wins}
\begin{itemize}
\item You can name the four pipeline stages and what each produces.
\item You can read \texttt{objdump -h/-d/-t/-r} output and find code,
data, symbols and relocations.
\item You know a \emph{link} error from a \emph{compile} error and how
to fix each.
\item You can place Book 5's cpu8 inside a real address space and name
the piece of hardware that made it real: the page table.
\end{itemize}
\end{wins}

\section{Common mistakes}
\label{sec:mistakes01}

\begin{itemize}
  \item \textbf{``Compile'' is one step.} It is four. Students lose hours
        calling a linker error a compiler bug; the fix is always
        different.
  \item \textbf{\texttt{.o} is ``half a program''.} It is a whole ELF,
        just not yet executable --- sections, symbols, relocations are
        already real. Nothing is created at link time except the
        \emph{layout}.
  \item \textbf{\texttt{.bss} wastes disk.} It does not exist in the
        file; the loader zero-fills pages. A gigabyte of zeros costs
        zero bytes on disk.
  \item \textbf{Disassembly ``proves'' optimisation.} \texttt{objdump -d}
        shows bytes; -O0 vs -O2 differences are layout and constants,
        not magic. Measure before attributing speed.
  \item \textbf{``The CPU runs the file.''} The CPU runs \emph{mapped
        pages}; the file is only the blueprint. \texttt{execve} is the
        step that turns one into the other.
\end{itemize}

\section*{Exercises}
\label{sec:exercises01}
\addcontentsline{toc}{section}{Exercises}

% ---- inline checkpoints (2, embedded above) already count; additional graded:
\begin{exercise}\label{ex:easy01}
On \texttt{ch01\_sum.o}, which section holds the machine code, and which
holds the relocation records that patch it?
\end{exercise}
\begin{solution}[hint]
\texttt{objdump -h} lists both by name: the code lives in one, the patch
list in another whose name starts with the code's.
\end{solution}

\begin{exercise}\label{ex:easy02}
\texttt{objdump -t} on \texttt{ch01\_sum.elf} lists \texttt{sum} as
\texttt{g F .text}. What do the two flags mean?
\end{exercise}
\begin{solution}[hint]
One flag says ``global'' (visible outside this file); the other says
``function''. Undefined entries would show \texttt{U} instead.
\end{solution}

\begin{exercise}\label{ex:easy03}
A program's \texttt{.bss} is 500~KB but its file on disk is 40~KB.
Why is this not a contradiction?
\end{exercise}
\begin{solution}[hint]
\texttt{.bss} records a size and a destination, not bytes; the loader
zero-fills the pages it names.
\end{solution}

\begin{exercise}\label{ex:easy04}
Your build prints \texttt{undefined symbol: main}
(from \texttt{ld.lld}). Which of the four pipeline stages failed?
\end{exercise}
\begin{solution}[hint]
The stage that resolves names --- the one that runs \emph{after}
\texttt{.o} files exist.
\end{solution}

\begin{exercise}\label{ex:med01}
\texttt{callq} in \texttt{main} encodes \texttt{00 00 00 00}. In one
sentence each: who left that hole, and who fills it?
\end{exercise}
\begin{solution}[hint]
The stage that did not know the address left it; the stage that learns
final addresses fills it via a relocation record.
\end{solution}

\begin{exercise}\label{ex:med02}
Compare \texttt{objdump -h} on \texttt{ch01\_sum.o} and
\texttt{ch01\_sum.elf}: what did the linker add, and what stayed
the same?
\end{exercise}
\begin{solution}[hint]
Addresses (VMA) appeared; section \emph{names} and \emph{jobs} did not
change. One section disappeared into runtime metadata.
\end{solution}

\begin{exercise}\label{ex:med03}
\texttt{execve} maps \texttt{.text} read-execute and \texttt{.data}
read-write. Which single bit of your cpu8's memory would have had to
become two bits, and why?
\end{exercise}
\begin{solution}[hint]
The page-table permission bits your flat RAM did not need ---
\texttt{.text} must not be writable, \texttt{.data} must not be
executable.
\end{solution}

\begin{exercise}\label{ex:med04}
\texttt{-O0} vs \texttt{-O2} on \texttt{ch01\_sum.c} changed
\texttt{.text} from 0x76 to 0x96 bytes. Was the faster version larger
or smaller here, and what does that tell you about ``optimise =
smaller''?
\end{exercise}
\begin{solution}[hint]
Faster was \emph{larger} (0x96 $>$ 0x76) --- unrolling/inlining trades
size for speed. Optimisation is about time, not bytes.
\end{solution}

\begin{exercise}[hard]\label{ex:hard01}
Hand-link the two objects from
Example~\ref{ex:handlink} and explain, per relocation record, which
hole each resolves and which it cannot. Then name the single symbol
whose absence forces the linker to fail.
\end{exercise}
\begin{solution}[hint]
The \texttt{PC32} record resolves within the same file; the two
\texttt{PLT32} records need other objects: the \texttt{sum} record is
satisfied once the second object joins the link, while the
\texttt{missing} record has no home anywhere --- that one is the error.
\end{solution}

\begin{exercise}[hard]\label{ex:hard02}
Design the minimum change to your cpu8 that would let it run an ELF
\texttt{.text} section directly. What two pieces of metadata does
\texttt{.text} carry that your ROM already has, and what two does
your ROM lack?
\end{exercise}
\begin{solution}[hint]
Your ROM already has bytes and an implied base address. It lacks a
\texttt{VMA} (which virtual page it lands on) and a protection mode
(read-execute vs read-write) --- the two fields that make an ELF a
program rather than a tape.
\end{solution}

\begin{exercise}[hard]\label{ex:hard03}
A student claims the entry point \texttt{\_start} is ``where the OS
starts the program''. Correct the claim in one sentence, then name
the field in the ELF header that actually points to the first
instruction.
\end{exercise}
\begin{solution}[hint]
\texttt{\_start} is not magic --- it is the symbol the linker recorded
in the ELF header's \texttt{e\_entry} field; the loader reads that
address and jumps to it, whoever named the symbol.
\end{solution}

\section*{Summary}
\label{sec:summary01}
\addcontentsline{toc}{section}{Summary}

\begin{center}
\begin{tabularx}{\linewidth}{lX}
\toprule
\textbf{Stage} & \textbf{What it produces} \\
\midrule
Preprocess & expanded C (one translation unit) \\
Compile & assembly (\texttt{.s}) \\
Assemble & ELF object (\texttt{.o}): sections + symbols + relocations \\
Link & executable ELF: resolved symbols, assigned VMAs \\
Load (\texttt{execve}) & mapped pages, \texttt{pc} $\to$ entry VA \\
\midrule
\textbf{You can now} & \\
\midrule
read \texttt{objdump} & code, data, symbols, relocations at a glance \\
name the failure & compile vs assemble vs link vs load --- different
fixes \\
place cpu8 in the map & page table is the piece that makes it real \\
\bottomrule
\end{tabularx}
\end{center}

\section*{Further reading}
\label{sec:fr01}
\addcontentsline{toc}{section}{Further reading}

\begin{itemize}
  \item Randal E. Bryant and David R. O'Hallaron, \emph{Computer Systems:
        A Programmer's Perspective}, ch.\,7 ``Linking'' --- the standard
        walk-through of symbols and relocation this chapter compresses.
  \item Ulrich Drepper, ``How To Write Shared Libraries'', \S1
        (ELF basics, sections vs segments).
        \url{https://www.akkadia.org/drepper/dsohowto.pdf}
  \item Compiler Explorer documentation --- try \texttt{clang -S} on your
        own sum online and match the output to Example~\ref{ex:pipeline}.
        \url{https://github.com/compiler-explorer/compiler-explorer}
  \item The ELF specification (System V ABI, \S4--5) --- the byte-level
        reference for every field \texttt{objdump} printed.
  \item Linux \texttt{execve(2)} and \texttt{elf(5)} man pages ---
        the loader's side of the contract.
\end{itemize}

\section*{Watch instead}
\label{sec:watch01}
\addcontentsline{toc}{section}{Watch instead}

\begin{rulebox}
\textbf{Watch instead.} If a video explains it faster than the prose,
watch these and come back:
\begin{itemize}
  \item ``In-depth: ELF --- The Extensible \& Linkable Format'' ---
        stacksmashing, 19:02. Byte-level walk of the format this
        chapter uses; matches Sections~\ref{sec:elf01}--\ref{sec:sym01}.
        \url{https://youtu.be/nC1U1LJQL8o}
  \item ``Writing a Linux ELF Binary by Hand'' --- Simon Racz, 59:11.
        Builds the file from hex --- the same journey with no toolchain.
        \url{https://youtu.be/JM9jX2aqkog}
\end{itemize}
\end{rulebox}

\chapternotes
